\documentclass[../../../include/open-logic-section]{subfiles}

\begin{document}

\olfileid{sth}{story}{blv}

\olsection{ضمیمہ۔ فریگے دا پنجواں بنیادی قانون}

\olref[rus]{sec} وچ اسیں سمجھایا سی کہ رسل نے اپنا پیراڈوکس اوس نظام لئی
اک مسئلے دے طور اُتے پیش کیتا جیہڑا فریگے نے اپنی \emph{Grundgesetze} وچ
بیان کیتا۔ فریگے دے نظام وچ سادہ خاصیت توں سیٹ بناؤن دا اصول سیدھی صورت
وچ شامل نہیں سی۔ سو ایس ضمیمے وچ اسیں بہت مختصر سمجھاواں گے کہ فریگے دے
نظام وچ \emph{کیہ شامل سی}، اوہدا سادے خاصیت توں سیٹ بناؤن دے اصول نال
کیہ تعلق اے تے رسل دے پیراڈوکس نال کیہ تعلق اے۔

فریگے دا نظام \emph{دوجے درجے دا} اے، تے ایس نوں \emph{کسے تصور دے
دائرے} دی دھارنا ٹھیک طرح لکھن لئی بنایا گیا سی۔\footnote{ٹھیک گل ایہ اے
کہ فریگے اک ہور عام دھارنا رسمی طور اُتے لکھن دی کوشش کردا اے، فنکشن
دی ’’قدراں دا دائرہ‘‘۔ تصوراں دے دائرے اوس عام دھارنا دی اک خاص صورت
نیں۔ تفصیل لئی \citet[pp.\ 8--9]{Heck2012} ویکھو۔} فریگے توں ملی علامت
ورت کے اسیں $\fregeext{x}{F(x)}$ نوں \emph{تصور~$F$ دا دائرہ} لکھاں گے۔
ایہ اک ترکیب اے جیہڑی \emph{محمول} ’’$F$‘‘ لے کے اینوں (پہلے درجے دی)
\emph{ٹرم} ’’$\fregeext{x}{F(x)}$‘‘ بنا دیندی اے۔ ایس ترکیب نال فریگے
نے رکنیت دی ہِیٹھلی \emph{تعریف} دتی:
\[
	a \in b =_\text{df} \exists G(b = \fregeext{x}{G(x)} \land Ga)
\]
موٹے طور اُتے، $a \in b$ اُدوں تے صرف اُدوں جدوں $a$ ایسے تصور تھلّے آؤندا
ہووے جیہدا دائرہ $b$ اے۔ (دھیان دیو کہ مقداری نشان ’’$\exists G$‘‘ دوجے
درجے دا اے۔) فریگے نے ہِیٹھلا اصول وی منّیا، جیہنوں \emph{پنجواں بنیادی قانون}
آکھدے نیں:
$$\fregeext{x}{F(x)} = \fregeext{x}{G(x)} \liff \forall x (Fx \liff Gx)$$
موٹے طور اُتے، تصوراں دے دائرے اکّو اُدوں تے صرف اُدوں نیں جدوں تصور
عین اوہناں ای شےواں اُتے لاگو ہوندے نیں۔ (فیر، ’’$F$‘‘ تے ’’$G$‘‘
دونیں محمول دی تھاں نیں۔) ہُن اک سدھا اصول رکنیت نوں خاصیت پوری کرن نال
جوڑدا اے:

\begin{lem}[\emph{Grundgesetze} وچ]\ollabel{lem:Fregeextensions}
$\forall F \forall a(a \in \fregeext{x}{F(x)} \liff Fa)$
\end{lem}

\begin{proof}
$F$ تے $a$ مقرر کرو۔ ہُن $a \in \fregeext{x}{F(x)}$ اُدوں تے صرف اُدوں
جدوں $\exists G(\fregeext{x}{F(x)} = \fregeext{x}{G(x)} \land Ga)$
(رکنیت دی تعریف نال)، اُدوں تے صرف اُدوں جدوں
$\exists G(\forall x(Fx \liff Gx) \land Ga)$ (پنجویں بنیادی قانون نال)،
اُدوں تے صرف اُدوں جدوں $Fa$ (دوجے درجے دی بنیادی منطق نال)۔
\end{proof}

تے ایس توں سادہ خاصیت توں سیٹ بناؤن دا اصول لگ بھگ سدھا نکلدا اے:

\begin{lem}[\emph{Grundgesetze} وچ۔]
$\forall F \exists s \forall a (a \in s \liff Fa)$
\end{lem}

\begin{proof}
$F$ مقرر کرو۔ ہُن \olref{lem:Fregeextensions} توں
$\forall a (a \in \fregeext{x}{F(x)} \liff Fa)$ ملدا اے۔ سو وجودی
تعمیم نال $\exists s\forall a(a \in s \liff Fa)$۔ نتیجہ ایس لئی نکلدا اے
کہ $F$ کوئی وی ہو سکدا سی۔
\end{proof}

$F$ نوں $\forall x(Fx \liff x \notin x)$ مطابق لین نال رسل دا پیراڈوکس نکلدا اے۔

\end{document}
